% Section 2: Discoveries and types of radiation.

\cstopic{2.1}{Discoveries from 1896 to 1911}{L4 pp.1--15}
\takeaway{This topic covers the four discoveries the course treats as the turning points: radioactivity itself, new radioactive elements, the three kinds of radiation and the nucleus. Exam questions ask who, when, what the experiment showed, and why it mattered.}

\begin{tablebox}{Timeline}{@{}>{\raggedright\arraybackslash}p{0.09\linewidth}>{\raggedright\arraybackslash}p{0.24\linewidth}X@{}}
\toprule\textbf{Year} & \textbf{Who} & \textbf{What it showed}\\\midrule
1896 & Henri Becquerel & uranium salts darken a wrapped photographic plate with no light: radiation comes from the atom itself\\
1898 & Marie and Pierre Curie & polonium and radium; radioactivity is a property of the element, not of a compound\\
1899 & Ernest Rutherford & two kinds of rays, alpha and beta, told apart by how far they penetrate\\
1911 & Rutherford, Geiger, Marsden & the gold-foil result: a small, dense, positive nucleus\\
\bottomrule
\end{tablebox}

\begin{factcard}{The gold-foil experiment}[figs/gold-foil.pdf][0.42]
\fact{Who, when}{Geiger and Marsden, under Rutherford, 1909 to 1911}
\fact{Setup}{alpha particles fired at gold foil a few hundred atoms thick}
\fact{Result}{almost all pass straight through; about 1 in 8{,}000 bounces back}
\fact{Why it matters}{the positive charge and nearly all the mass sit in a tiny nucleus; the old ``plum pudding'' model could not turn an alpha particle back}
\end{factcard}

\begin{trap}{Discovered versus explained}
Becquerel discovered radioactivity in 1896, but the Curies named it and showed it belongs to the element. Questions that ask who ``showed radioactivity is atomic'' want the Curies.
\end{trap}

\cstopic{2.2}{Alpha, beta and gamma radiation}{L5 pp.2--10}
\takeaway{This topic covers the three kinds of radiation, what each is made of, what it does to the nucleus that emits it, and how far it travels. The pattern to remember: the heavier and more charged the particle, the less it penetrates.}

\begin{tablebox}{The three kinds side by side}{@{}>{\raggedright\arraybackslash}p{0.13\linewidth}>{\raggedright\arraybackslash}p{0.25\linewidth}>{\raggedright\arraybackslash}p{0.27\linewidth}X@{}}
\toprule & \textbf{What it is} & \textbf{Change to the nucleus} & \textbf{Stopped by}\\\midrule
alpha & helium nucleus: 2 protons, 2 neutrons & mass number \(-4\), atomic number \(-2\) & a sheet of paper or skin\\
beta & electron from a neutron turning into a proton & mass number same, atomic number \(+1\) & a few millimetres of aluminium\\
gamma & high-energy photon & no change in either number & several centimetres of lead\\
\bottomrule
\end{tablebox}

\begin{keybox}{Words the exam uses}
\defn{Isotope}{atoms of the same element (same number of protons) with different numbers of neutrons.}
\defn{Ionizing}{able to knock electrons off atoms; alpha ionizes most per centimetre, gamma least.}
\defn{Decay chain}{a sequence of decays, each daughter itself unstable, ending at a stable nucleus.}
\end{keybox}

\begin{tfbox}
\tf{T}{Alpha radiation is the most ionizing but the least penetrating.}{Heavy and doubly charged, so it loses energy fast.}
\tf{F}{Gamma decay changes the element.}{It only lowers the nucleus's energy.}
\tf{T}{In beta decay the atomic number goes up by one.}{A neutron becomes a proton.}
\end{tfbox}
